\subsection{Multi-Seed Scalability Results}
\label{sec:multi_condition}
\label{sec:multi_seed}
\label{sec:dqn_results}
\label{sec:scalability}

Results of the five-condition scalability sweep (50 seeds per condition, 250 episodes per agent) follow. Full numerical results are in Table~\ref{tab:scalability_ndr} (Section~\ref{sec:relational_headtohead}) and Figure~\ref{fig:scalability_combined}.

At $(100{\times}100, N{=}10)$ all agents reach $100\%$ NDR. From $N{=}20$ upward the DQN drops sharply: $99.4\%$, $52.7\%$, $27.5\%$, $13.0\%$ across the four larger conditions, while the TSP Oracle stays at $76.8\%$ at $N{=}50$, a $63.8$ percentage-point gap that signals a routing bottleneck, not a physical-link limit.

The greedy baselines beat the DQN from $N\ge 30$. SF-Aware Greedy leads $84.3\%$ vs $52.7\%$ at $N{=}30$ ($p < 10^{-4}$), and Nearest-Sensor Greedy, which has no SF awareness at all, also overtakes the DQN from $N{=}30$ onward. A protocol-blind heuristic beating a protocol-aware learner points to the routing representation as the problem, not the protocol stack. Section~\ref{sec:diagnostic} traces this to the flat-MLP encoder's inability to treat the sensor set as an unordered collection.

Jain's Fairness Index tracks the same pattern as NDR. At $(100{\times}100, N{=}10)$ all agents achieve high fairness ($\ge 0.78$), though the DQN already trails both greedy agents ($0.78$ vs $1.00$) because its perimeter trajectory collects repeatedly from the same sensors while interior ones go unserved. The gap grows with scale: at $(500{\times}500, N{=}50)$ the DQN's index falls to $0.08$, meaning nearly all collected data comes from the handful of sensors the perimeter path passes near. The TSP Oracle reaches only $0.46$ at the same condition, so the battery-limited episode sets a fairness ceiling well below $1.00$ even with optimal routing. Full results are in Table~\ref{tab:scalability_ndr} (Section~\ref{sec:relational_headtohead}).
